% ==============================================================================
% CAPÍTULO 5: CONCLUSÃO (somente na Defesa)
% Substitui automaticamente o capítulo "Planos para a continuidade do trabalho"
% quando \Tipoexame{Defesa} (ver main.tex). Também pode ser intitulado
% "Considerações finais".
% A conclusão NÃO é um resumo do trabalho (manual, cap. 3, item 3): é a síntese
% dos resultados à luz da pergunta de pesquisa e dos objetivos.
% Não se introduzem dados, citações ou argumentos novos neste capítulo.
% ==============================================================================
\chapter{Conclusão}
\label{ch:conclusao}

\preencher{A conclusão tem como finalidade responder explicitamente à pergunta de pesquisa e demonstrar em que medida os objetivos foram atingidos, com base nos resultados apresentados no Capítulo~\ref{ch:analise_discussao}. Apagar este texto de instrução após a redação do capítulo.}

\section{Resposta à pergunta de pesquisa}
\label{sec:resposta_pergunta}
Retomar a \hyperref[perg:principal]{Pergunta de Pesquisa} e respondê-la de forma direta, indicando como cada objetivo específico (Seção~\ref{subsec:objetivos_especificos}) foi atendido e com quais evidências. Convém seguir a mesma ordem dos objetivos, o que permite ao leitor verificar a correspondência entre o que foi proposto e o que foi alcançado.

\section{Contribuições teóricas e práticas}
\label{sec:contribuicoes_finais}
Explicitar o que o trabalho acrescenta à literatura e à prática, confrontando o resultado obtido com as contribuições esperadas anunciadas na Seção~\ref{sec:contribuicoes}. Contribuições não confirmadas também devem ser reconhecidas.

\section{Limitações}
\label{sec:limitacoes_finais}
Apresentar as limitações efetivamente verificadas ao longo da pesquisa, retomando as previstas na Seção~\ref{sec:limitacoes} e indicando de que modo elas restringem o alcance das conclusões.

\section{Trabalhos futuros}
\label{sec:trabalhos_futuros}
Indicar desdobramentos concretos, preferencialmente derivados das limitações identificadas, e não sugestões genéricas.
